\documentclass[final]{beamer}

\usepackage[orientation=portrait,size=a0,scale=1.6]{beamerposter}
\usepackage[utf8]{inputenc}
\usepackage[T1]{fontenc}
\usepackage[ngerman]{babel}
\usepackage{xcolor}
\usepackage{tikz}
\usepackage{lmodern}
\usepackage{graphicx}
\usepackage{pifont}

\definecolor{primary}{HTML}{0D9488}
\definecolor{dark}{HTML}{134E4A}
\definecolor{muted}{HTML}{6B7280}
\definecolor{soft}{HTML}{F0FDFA}
\definecolor{warn}{HTML}{D97706}

\usetheme{default}
\setbeamercolor{background canvas}{bg=white}
\setbeamercolor{block title}{fg=white,bg=primary}
\setbeamercolor{block body}{fg=dark,bg=soft}
\setbeamercolor{block title alerted}{fg=white,bg=dark}
\setbeamercolor{block body alerted}{fg=dark,bg=soft}
\setbeamerfont{block title}{size=\LARGE,series=\bfseries}
\setbeamerfont{block body}{size=\large}
\setbeamertemplate{navigation symbols}{}

\newlength{\sepwidth}\setlength{\sepwidth}{0.025\textwidth}
\newlength{\boxwidth}\setlength{\boxwidth}{0.305\textwidth}

\begin{document}
\begin{frame}[t]{}

% ── HEADER ───────────────────────────────────────────────────────────────────
\begin{tikzpicture}[remember picture,overlay]
  \fill[primary] (current page.north west)
    rectangle ([yshift=-24cm]current page.north east);
\end{tikzpicture}

\vspace{1.0cm}
{\color{white}\fontsize{52}{62}\selectfont\bfseries
  Wie komplex ist eine Website wirklich?}\\[0.6cm]
{\color{white}\fontsize{24}{30}\selectfont
  Jonas Kimmer · Wirtschaftsinformatik M.Sc. · Hochschule RheinMain · 2026}\\[0.3cm]
{\color{soft}\fontsize{20}{26}\selectfont
  Betreuung: M.Sc. Yasmina Tajja \quad|\quad j.kimmer@student.hs-rm.de}
\vspace{2.0cm}

% ── ZEILE 1: Forschungsfrage + Beitrag ───────────────────────────────────────
\begin{columns}[T]
\begin{column}{0.47\textwidth}
  \begin{alertblock}{Forschungsfrage}
    \vspace{0.3em}
    \large
    Website-Komplexität lässt sich heute nicht\\
    objektiv messen oder vergleichen.\\[0.5em]
    \LARGE\bfseries
    Lässt sich Website-Komplexität automatisch\\
    messen und klassifizieren, ohne\\
    manuelle Inspektion?
    \vspace{0.3em}
  \end{alertblock}
\end{column}
\hspace{\sepwidth}
\begin{column}{0.47\textwidth}
  \begin{block}{Beitrag dieser Arbeit}
    \vspace{0.4em}
    {\large
    \textcolor{primary}{\bfseries Tool\enskip}--\enskip
      Python · Streamlit · Einlesen · 6 Merkmale · Export\\[0.5em]
    \textcolor{primary}{\bfseries Daten\enskip}--\enskip
      DOM-Tiefe · Knoten · Links · Text · Medien · Formulare\\[0.5em]
    \textcolor{primary}{\bfseries Ergebnis\enskip}--\enskip
      Clustering ohne manuelle Labels
    }
    \vspace{0.3em}
  \end{block}
\end{column}
\end{columns}

\vspace{0.8em}

% ── ZEILE 2: Cluster-Bild (grüner Balken, Titel zentriert) ───────────────────
{\setbeamertemplate{block title}{\vskip0.4em\centering\LARGE\bfseries\insertblocktitle\vskip0.4em}
\begin{block}{Seitentypen automatisch erkannt}
  \vspace{0.4em}
  \begin{center}
    \includegraphics[width=0.62\linewidth]{paper2_abb2_cluster_scatter.png}\\[0.4em]
    {\large\color{muted}
    Proof of Concept · 3 Test-Websites · 47 Seiten · drei Typen vollautomatisch gefunden}
  \end{center}
  \vspace{0.2em}
\end{block}}

\vspace{0.8em}

% ── ZEILE 3: 3 Befunde nebeneinander ─────────────────────────────────────────
\begin{columns}[T]

\begin{column}{\boxwidth}
  \begin{tikzpicture}
    \node[fill=primary,rounded corners=8pt,inner sep=14pt,
          text width=0.95\linewidth]{
      \color{white!70!primary}\large\bfseries
      \ding{172} Zwei Merkmale reichen\\[0.5em]
      \color{white}\large\mdseries
      Linkanzahl und DOM-Tiefe allein\\
      genügen zur Klassifikation.\\
      Alle anderen Merkmale sind zweitrangig.
    };
  \end{tikzpicture}
\end{column}

\hspace{\sepwidth}

\begin{column}{\boxwidth}
  \begin{tikzpicture}
    \node[fill=dark,rounded corners=8pt,inner sep=14pt,
          text width=0.95\linewidth]{
      \color{white!60!dark}\large\bfseries
      \ding{173} Kein einziges Label von Hand\\[0.5em]
      \color{white}\large\mdseries
      Normales ML braucht manuell\\
      markierte Beispiele. Hier klassifiziert\\
      der Algorithmus vollständig selbst.
    };
  \end{tikzpicture}
\end{column}

\hspace{\sepwidth}

\begin{column}{\boxwidth}
  \begin{tikzpicture}
    \node[fill=warn,rounded corners=8pt,inner sep=14pt,
          text width=0.95\linewidth]{
      \color{white!70!warn}\large\bfseries
      \ding{174} Formulare? Fast egal.\\[0.5em]
      \color{white}\large\mdseries
      Obwohl Formulare die stärkste\\
      Interaktion bieten, erklären sie\\
      Komplexität am wenigsten.
    };
  \end{tikzpicture}
\end{column}

\end{columns}

\vspace{0.8em}

% ── ZEILE 4: Ausblick (volle Breite) ─────────────────────────────────────────
\begin{alertblock}{Anwendung \& Ausblick}
  \vspace{0.3em}
  \LARGE\bfseries
  Komplexität ist messbar. Jetzt muss der Mensch bestätigen.\\[0.5em]
  \large\mdseries
  · Einsetzbar für Redesign-Vergleiche, Website-Audits und automatische Kategorisierung\\[0.3em]
  · Offene Frage: Stimmen Algorithmus und Mensch überein?\\[0.3em]
  {\color{muted!60!white}\normalsize Manuelle Validierung auf einem größeren Datensatz steht aus.}
  \vspace{0.3em}
\end{alertblock}

\end{frame}
\end{document}
